\documentclass[10pt,a4paper]{amsart}
\usepackage[margin=19mm]{geometry}
\usepackage{amsmath,amssymb,mathtools}
\usepackage[scr=rsfs,cal=euler]{mathalfa}
\usepackage{xfrac}
\usepackage[unicode,hidelinks,bookmarks=false]{hyperref}
\newcommand{\ie}{i.e.\ }
\newcommand{\cf}{cf.\ }
\newcommand{\resp}{respectively\ }
\newcommand{\Subsection}{\subsection}
\providecommand{\nobreakdash}{\nobreak\hbox{-}}
\setlength{\emergencystretch}{2em}
\multlinegap0pt
\usepackage{bm}
\usepackage{bbm}
\usepackage{dsfont}
\usepackage{xspace}
\usepackage{cancel}
\usepackage{mathtools}
\usepackage{amsthm}
\makeatletter
\@ifundefined{cla}{\newtheorem{cla}{Claim}}{}
\makeatother
\newcommand\cB{{\mathcal B}}
\newcommand\cC{{\mathcal C}}
\newcommand\cE{{\mathcal E}}
\newcommand\cH{{\mathcal H}}
\title[The Tur\'{a}n number of Berge matchings]{The Tur\'{a}n number of Berge matchings}
\author{Yichen Wang, Zixuan Yang, Xiamiao Zhao, Yuhang Bai, Junpeng Zhou}
\date{Source: arXiv:2510.05422v2 (CC BY 4.0)}
\begin{document}
\maketitle

\noindent\emph{Source and licence.} The Tur\'{a}n number of Berge matchings. Version arXiv:2510.05422v2 (CC BY 4.0), distributed under the Creative Commons Attribution 4.0 International licence, as recorded in the arXiv licence metadata for this exact version and shown on the abstract page https://arxiv.org/abs/2510.05422v2. Full rights notice: licenses/arxiv-2510.05422-CC-BY-4.0.txt.\par\smallskip

\noindent\textit{Editorial scope.} Counted unit: the Cla whose source label is \texttt{claim: 2 isolated vtx} in the article (source0.tex lines 598--602, source label \texttt{claim: 2 isolated vtx}), delivered together with the article's own complete original proof of it (source0.tex lines 604--638, one \texttt{proof} environment of 35 lines). No mathematics is written, repaired or re-derived: statement and proof are the article's own text line for line. Required context retained uncounted: none. Every definition, display and label of those lines is retained unchanged. Omitted from this package and not redistributed: 1--597, 603--603, 639--1413. Nothing inside a retained excerpt is omitted or altered: every retained body is delivered exactly as the article's lines are written, so no omission receipt of any kind accompanies this package. Citations: the retained excerpts carry no citation command at all, so no bibliography entry of the article is transcribed and no reference is redistributed. Reference adaptation: none was needed and none was made; every reference inside the delivered unit resolves inside it, so no unresolved reference or citation is produced. Printed slips kept literal: no printed word, symbol, hypothesis or number of the counted statement or of its proof was altered. Layout and class: the article's own class is \texttt{article} with packages including amsmath, amssymb, authblk, geometry, hyperref, graphicx, epstopdf, mathrsfs, amsfonts, amscd, amsthm, bm, psfrag, natbib; the wrapper is the lane's pinned \texttt{amsart} editorial wrapper at 10pt on A4, which restates the theorem-like environments the retained text needs and adds the article's own macro definitions that the retained text needs (\texttt{\textbackslash cB}, \texttt{\textbackslash cC}, \texttt{\textbackslash cE}, \texttt{\textbackslash cH}), with extra wrapper packages (bm, bbm, dsfont, xspace, cancel, mathtools). The delivered numbering is the unit's own. Assets: no retained span contains a figure environment, a graphics-inclusion command, a PDF-page inclusion, a table or a TikZ picture; no image file is copied into this package, the package declares no asset dependency and no image file is redistributed with it. Licence: version arXiv:2510.05422v2 of this article is distributed under the Creative Commons Attribution 4.0 International licence, as recorded in the arXiv licence metadata for this exact version and shown on the abstract page https://arxiv.org/abs/2510.05422v2. A full rights notice is delivered at licenses/arxiv-2510.05422-CC-BY-4.0.txt. Characters: every retained excerpt is pure ASCII, so the delivered engine, XeLaTeX with its default Latin Modern text font, reports no missing character.
    \begin{cla}\label{claim: 2 isolated vtx}
        If $s\leq r\leq 2s$,
        then either $e(\cH)\leq \binom{2s+1}{r}$,
        or there are at least two vertices $u, v\in V(\cH)\setminus \cC$ such that $d_{\cH'}(u) \geq 1$ and $d_{\cH'}(v) \geq 1$.
    \end{cla}

    \begin{proof}[{\bf{Proof of Claim~\ref{claim: 2 isolated vtx}.}}]
    %\noindent\textbf{Proof of Claim~\ref{claim: 2 isolated vtx}:}
    It is sufficient to prove that if there is at most one vertex $u\in V(\cH)\setminus \cC$ with $d_{\cH'}(u) \geq 1$, then $e(\cH)\leq \binom{2s+1}{r}$.
    If $d_{\cH'}(v)=0$ for every $v\in V(\cH)\setminus \cC$, then $e(\cH) \leq s+\binom{2s}{r}\leq \binom{2s+1}{r}$ as $r\leq 2s$.

    Let $v\in V(\cH)\setminus \cC$ be the unique vertex with $d_{\cH'}(v)\geq 1$. If $e_i \subseteq \cC \cup \{v\}$ for every $i\in [s]$, then $e(\cH) \leq \binom{2s+1}{r}$ and we are done. Now suppose that there exists some $i\in [s]$ such that $e_i \nsubseteq \cC \cup \{v\}$.
    Without loss of generality, suppose $e_s$ contains a vertex $w \notin \cC \cup \{v\}$. Note that $|\cC|=2s$. If $r=2s$, then we have $u_s\in N_{\cH'}(v)$ or $v_s\in N_{\cH'}(v)$. Thus, there exists a $\cB M_{s+1}$ in $\cH$ by adding the hyperedge containing $v$ and $u_s$~(or $v_s$) to $\cE$, and extending the core to $\cC\cup \{v,w\}$, a contradiction. Now we assume that $s\leq r<2s$.
    %Note that $N_{\cH'}(v) \subseteq \cC$. Indeed, if there exists a vertex $u\in V(\cH)\setminus \cC$ such that $\{u,v\}\subseteq e\in E(\cH')$, then $\cE\cup \{e\}$ is a Berge copy of $M_{s+1}$  with the core $\cC\cup \{v,u\}$ of $\cH$, a contradiction.
    %Next we consider two cases: one is $e_i \subseteq \cC \cup \{v\}$ for every $i\in [s]$, and the other is that there exists $e_i \nsubseteq \cC \cup \{v\}$ for some $i\in [s]$.
    %For the first case, we have  $e(\cH) \leq \binom{2s+1}{r}$.
    We claim that $u_s\notin N_{\cH'}(v)$ and $v_s \notin N_{\cH'}(v)$.
    Indeed, if $u_s\in N_{\cH'}(v)$ (or $v_s\in N_{\cH'}(v)$), then as in the case $r=2s$ we may find a $\cB M_{s+1}$ in $\cH$, a contradiction.
    %If $r=2s$, then we have $u_s\in N_{\cH'}(v)$ or $v_s\in N_{\cH'}(v)$. So there exists a $\cB M_{s+1}$ in $\cH$, a contradiction.

    Then $|E_{\cH'}(v)|\leq \binom{|N_{\cH'}(v)|}{r-1}$.
    Note that $E(\cH'[\cC])=E(\cH')\setminus E_{\cH'}(v)$.
    %Note that every hyperedge of $E(\cH')\setminus E_{\cH'}(v)$ is contained in the vertex set $\cC$.
    We now partition $E(\cH'[\cC])$ into the following three classes:
    $\cE_1$ (consisting of hyperedges not containing $\overline{N_{\cH'}(v)}$),
    $\cE_2$ (consisting of hyperedges not containing $\{u_s,v_s\}$),
    and $\cE_3$ (consisting of hyperedges containing both a vertex of $\{u_s,v_s\}$ and a vertex of $\overline{N_{\cH'}(v)}$).
    %We now consider  the edges of $E(\cH')\setminus E_{\cH'}(v)$ formed by $\cE_1$ (hyperedges  not containing $\overline{N_{\cH'}(v)}$), $\cE_2$ (hyperedges not containing $\{u_s,v_s\}$),  and $\cE_3$ (hyperedges containing both  one vertex in $\{u_s,v_s\}$ and one  in $\overline{N_{\cH'}(v)}$).
    We assert that $\cE_3=\emptyset$. Otherwise, assume that $e\in \cE_3$ contains the vertex $u_s$ and $u_i\in \overline{N_{\cH'}(v)}$ (or $v_i\in \overline{N_{\cH'}(v)}$) for some $i\in [s-1]$.
    By the definition of $\overline{N_{\cH'}(v)}$, there is a hyperedge $e'\in E(\cH')$ containing $v$ and $v_i$ (or $u_i$). Then $(\cE\setminus \{e_i\})\cup \{e',e\}$ forms a $\cB M_{s+1}$ in $\cH$ with the core $\cC\cup \{v,w\}$, a contradiction.
    %i.e., for any hyperedge $e\in E(\cH')\setminus  E_{\cH'}(v)$, $e$ can not contain both one vertex in $\{u_s,v_s\}$ and one  in $\overline{N_{\cH'}(v)}$.
    %Suppose, without loss of generality, that $e\in \cE_3$ contains the vertex $u_s$ and $u_i\in \overline{N_{\cH'}(v)}$ for some $i\in [s-1]$.
    Note that $r-1\leq|N_{\cH'}(v)|\leq 2s-2$. Therefore, we have
     \begin{equation*}
    \begin{aligned}
        e(\cH)= s+|E_{\cH'}(v)|+|\cE_1|+|\cE_2| &\leq s + \binom{|N_{\cH'}(v)|}{r-1} + \binom{2s-|N_{{\cH'}}(v)|}{r} + \binom{2s-2}{r}
        %&\leq s + \binom{2s-2}{r-1} + \binom{2s-r+1}{r}+\binom{2s-2}{r} \\
        \le \binom{2s+1}{r}.
    \end{aligned}
     \end{equation*}
     \end{proof}

\end{document}
